%% kapa: exam=1 round=36 year=2025 period=1 form=A subject=civil_law subject_ko=민법 seq=28 number=28
%% source: Q-Net 감정평가사 기출문제 (cst00309) artlSeq=5231525 file=raw/1008/5231525/2025년 제36회 감정평가사 1차시험 1교시 문제지 원본.pdf
%% conversion: kordoc 4.24.51 -> tools/gen1.py (auto); review=proofread-sonnet; flags=none
\begin{question}{28}
甲, 乙, 丙은 X토지를 동일한 지분비율로 공유하고 있다. 이에 관한 설명으로 옳지 않은 것은? (다툼이 있으면 판례에 따름)
\begin{choices}
  \item 甲은 특별한 사정이 없는 한 자신의 지분을 자유롭게 처분할 수 있다.
  \item 甲과 乙이 丙과의 협의없이 X토지에 건물을 신축하여 임대하기로 결정하는 것도 관리 방법으로서 적법하다.
  \item 甲, 乙, 丙은 X토지를 3년 동안 분할하지 아니할 것을 약정할 수 있다.
  \item 甲이 乙 및 丙과의 협의없이 X토지 전부를 독점적으로 점유하여 사용하는 경우, 乙은 甲에게 X토지의 인도를 청구할 수 없다.
  \item 丙의 지분위에 원인무효의 저당권 등기가 마쳐진 경우, 甲은 X토지의 보존행위로서 저당권 등기의 말소를 청구할 수는 없다.
\end{choices}
\end{question}
\answer{2}
